\section{Discussion}
\label{sec:discussion}

\textbf{$\dq$ is the signal, and we found nothing to add to it.} Every model
takes most of its information from $\dq$, which already sums the effect of
every mechanism. Over the first 100 cycles it grows linearly, so the early
data contain no sign of the coming knee beyond its rate.

\textbf{The batch, not the model, sets the error.} With two training
batches, anything that separates b1 from b2 looks useful in
cross-validation and on the primary set: the temperature integral and charge
time of F3, the narrow voltage weights of PLS, the extra flexibility of the
GAM, the overhang rise. None of it holds on a new batch. Screening within
each batch (Section~\ref{sec:rq3}) catches part of this, but a model or
feature should finally be judged on a batch it has not seen.

\textbf{The test procedure changes the relation between $\dq$ and life.} In
u1 the \texttt{newstructure} cells live about twice as long as standard
cells charged with the same policy (median 809--997 vs.\ 446--492 cycles for
the three shared policies), and their $\dq$ is smaller still than that longer
life implies.

\textbf{Why b2 and the u1 standard cells live short.} The two short-lived
sets (median 481 and 451 cycles) are the two with 5-minute rests after the
fast charge and after discharge (Table~\ref{tab:batches}; for b2 also stated
in~\cite{severson2019}). The policies do not explain this alone. b2 has only
10-minute policies, but b1 cells with 10-minute policies still live longer
(median 700). 4.8C(80\,\%)-4.8C gives 636--870 cycles in b1 and 461--495 in
b2, and the three policies shared by the two u1 groups give 412--514 cycles
with long rests and 777--1186 with short ones. One policy, 6C(60\,\%)-3C,
barely differs (731--757 in b1, 713 in b2). The end-of-life label is the
same in every batch (first cycle below 0.88\,Ah).

A possible mechanism is temperature: the cell cools during the rests and
starts the fast charge 1.5--2\,\textdegree C colder (31.4--31.6 vs.\
32.8--34.4\,\textdegree C), which favours lithium plating. Within a set and
policy, where the procedure is identical, colder cells do live shorter:
$+6$\,\% life per \textdegree C at the start of the fast charge (95\,\%
interval 4--9\,\%; 112 cells, 36 policy groups). This slope, however,
accounts for only a 9--11\,\% shorter life across the observed temperature
gaps, against 35--55\,\% observed. Temperature contributes, but most of the
effect of the long rests is not explained by the data. Other procedure
differences, such as the CV cut-off, and an unrecorded change behind the
\texttt{newstructure} label cannot be ruled out.

\textbf{Uncertainty needs an out-of-range check.} The GP interval is
too narrow when the GP nearly interpolates the training cells, and
overconfident when the procedure changes (56\,\% coverage on the
\texttt{newstructure} cells). Large input $z$-scores flag some failures, but
not the \texttt{newstructure} cells, whose inputs were only moderately
unusual.

\subsection{Limitations}
\begin{itemize}
  \item \textbf{Small sample.} 41 training cells and about 40 cells per test
  set: most differences between models have bootstrap intervals that include
  zero. Many configurations were scored on the same test sets; fixing every choice on the training cells
  and the bootstrap limit, but do not remove, the optimism this causes.
  \item \textbf{Narrow scope.} One cell type (LFP/graphite, A123), one
  temperature (30\,\textdegree C), two training batches and only nine
  \texttt{newstructure} cells.
  \item \textbf{Limited search.} Nineteen candidates were screened; other
  features (e.g.\ from the charge curve or relaxation voltage) were not
  tried.
  \item \textbf{Unlabelled procedure changes.} What \texttt{newstructure}
  changes is inferred from the step records (rests, CV cut-off); other
  changes cannot be ruled out.
  \item \textbf{Two-cycle features.} $\dq$ uses only cycles 10 and 100, so
  noise in either passes directly into every feature.
\end{itemize}

\section{Conclusion}
\label{sec:conclusion}
The data-driven early prediction of~\cite{severson2019} was reproduced on a
cleaned pipeline, exactly for the one-feature model and within a few cycles
for the others (RQ1). Nineteen features derived from degradation mechanisms
were screened, and none added usable information beyond $\dq$ and the
paper's capacity features (RQ3). With 41 training cells, a Gaussian process
and an additive spline model were not significantly better than the
elastic net, within the training batches or on a new batch; the spline
model fails beyond the training range, and a random forest, which cannot
extrapolate at all, is worse everywhere (RQ2). About 40 cycles give a rough
prediction, 100 cycles a good one, and more bring no gain (RQ4). On an
unseen batch, cells tested like a training batch are predicted to within
9--14\,\%, while cells tested with a different procedure are over-predicted
by 17--62\,\% by every model; SHAP traces this to $\dq$ itself (RQ5). The
main obstacle to early life prediction on this dataset is the change of test
procedure between batches, not the choice of model or features.

\textbf{Future work.}
\begin{itemize}
  \item Train on more batches and procedures, including u1, and evaluate
  leave-one-batch-out.
  \item Give the model the procedure (rest length, CV cut-off) as an input,
  or correct $\dq$ for it.
  \item Fit $\dq$ over all cycles 10--100 instead of two, using its linear
  growth, to reduce noise.
  \item Combine the GP interval with an out-of-range check before trusting a
  prediction.
  \item Look for new information in signals not used here, such as the
  charge curve or rest voltage relaxation, and test on other chemistries and
  in closed-loop protocol optimisation~\cite{attia2020}.
\end{itemize}
